\chapter{Design Methodology}
\begin{abstract}

This document summarizes the design procedure for supersonic impulse turbine blades,
based on classical shock–expansion theory and the method of characteristics.
The accompanying program, OpenLaval, implements this procedure and supports
both symmetric and asymmetric blade geometries by allowing independent
Prandtl–Meyer angles for the upper and lower surfaces.

In an asymmetric blade, the upper and lower surfaces undergo different amounts
of expansion, resulting in different R$^\ast$ values, different characteristic
curves, and different arc segments. This chapter describes the complete design
workflow, including the asymmetric formulation now supported by the software.

\end{abstract}

\section{Overview}

The design of a supersonic impulse turbine blade begins with constructing
wall shapes that cancel the shocks generated at the leading edge.
Only a small number of physical parameters are required:

\begin{itemize}
  \item inlet Mach number $M_{\mathrm{in}}$
  \item outlet Mach number $M_{\mathrm{out}}$
  \item specific heat ratio $\gamma$
  \item inlet flow angle $\beta_{\mathrm{in}}$
  \item Prandtl–Meyer angles for the upper and lower surfaces
\end{itemize}

OpenLaval supports both symmetric blades
(using a single pair $(\nu_l, \nu_u)$)
and asymmetric blades
(using four independent angles):


\[
\nu_{l,\mathrm{lower}},\quad
\nu_{u,\mathrm{lower}},\quad
\nu_{l,\mathrm{upper}},\quad
\nu_{u,\mathrm{upper}}.
\]



These angles must satisfy the physical constraint


\[
\nu_{\min} \le \nu_l < \nu_u \le \nu_{\max},
\]


where $\nu_{\min}$ and $\nu_{\max}$ are the Prandtl–Meyer angles
corresponding to $M_{\mathrm{in}}$ and $M_{\mathrm{out}}$.

\section{Design Procedure}

The complete design workflow is as follows:

\begin{enumerate}
  \item Generate shock-cancelling characteristic curves for the upper and lower surfaces.
  \item Construct arc segments for each surface (R$^\ast$ differs between surfaces in the asymmetric case).
  \item Construct the upper-surface straight-line region based on the inlet flow angle and trim it to match the lower surface.
  \item Determine the blade height from the vertical separation of the two surfaces and assemble the full coordinates.
  \item Apply leading-edge trimming using the parameters \texttt{edge\_delta} and \texttt{edge\_offset}.
\end{enumerate}

In an asymmetric blade, steps 1–3 are performed independently for the upper and lower surfaces.

\section{Asymmetric Prandtl–Meyer Angles}

In the asymmetric formulation, the designer specifies four angles:


\[
(\nu_{l,\mathrm{lower}}, \nu_{u,\mathrm{lower}}), \qquad
(\nu_{l,\mathrm{upper}}, \nu_{u,\mathrm{upper}}).
\]



These define the amount of expansion on each surface.
Because the characteristic curves depend directly on these angles,
the resulting geometry differs between the upper and lower sides.

OpenLaval validates all four angles to ensure they lie within the physically
admissible Prandtl–Meyer range.

\section{Shock-Cancelling Curves}

The shock-cancelling curves are generated from the characteristic relation


\[
\phi(R^\ast, \theta) = \mathrm{const.}
\]



In the asymmetric case, the initial conditions differ between the upper and
lower surfaces, producing two distinct characteristic curves and two distinct
R$^\ast$ values.

\section{Arc Segments}

To ensure continuity between the characteristic curves and the straight-line
region, each surface receives its own circular arc segment.
Because R$^\ast$ differs between surfaces in the asymmetric case,
the arc radii also differ.

\section{Straight-Line Region}

The upper-surface straight-line region is aligned with the inlet flow angle
$\beta_{\mathrm{in}}$.
It is then trimmed to match the lower-surface $x$-coordinate.
In the asymmetric case, the connection points differ between surfaces.

\section{Leading-Edge Trimming}

OpenLaval applies leading-edge trimming using the parameters
\texttt{edge\_delta} and \texttt{edge\_offset}.
This operation is performed after both surfaces have been constructed.
Because the surfaces differ in the asymmetric case, the trimming points also differ.

\section{Program Usage}

The configuration file supports the following parameters:

\begin{itemize}
  \item symmetric PM angles: \texttt{vl}, \texttt{vu}
  \item asymmetric PM angles:
        \texttt{vl\_lower}, \texttt{vl\_upper},
        \texttt{vu\_lower}, \texttt{vu\_upper}
  \item inlet flow angle: \texttt{beta\_in}
  \item leading-edge parameters: \texttt{edge\_delta}, \texttt{edge\_offset}
  \item interpolation resolution: \texttt{num\_points}
\end{itemize}

The CLI provides an \texttt{--asymmetric} flag and override options for all
asymmetric angles.
